% !TeX program = pdflatex
\documentclass[11pt]{article}

% -----------------------
% Packages (barebones)
% -----------------------
\usepackage[margin=1in]{geometry}
\usepackage{setspace}
\usepackage{microtype}
\usepackage{amsmath, amssymb, amsthm, mathtools}
\usepackage{bm}              % bold math symbols
\usepackage{graphicx}
\usepackage{booktabs}
\usepackage{hyperref}
\usepackage[backend=biber]{biblatex}
\addbibresource{Proposal_bib.bib}

\onehalfspacing
\hypersetup{
  colorlinks=true,
  linkcolor=blue,
  citecolor=blue,
  urlcolor=blue
}

% -----------------------
% Metadata macros
% -----------------------
\newcommand{\Title}{Returns to Representation: Girls' Youth Sports Participation Following the Founding of the WNBA}
\newcommand{\Author}{Tate Mason}
\newcommand{\Date}{\today}

% -----------------------
% Section helpers
% -----------------------
% Use: \Section{Motivation} ... \Section{Methods}
\newcommand{\Section}[1]{\section{#1}}
\newcommand{\Subsection}[1]{\subsection{#1}}
\newcommand{\Subsubsection}[1]{\subsubsection{#1}}

% -----------------------
% Equation helpers
% -----------------------
% Inline: \(\E[Y|X]\), display: \eqn{ ... }, numbered: \eqnlabel{key}{...}
\newcommand{\eqn}[1]{\begin{equation*}#1\end{equation*}}
\newcommand{\eqnlabel}[2]{\begin{equation}\label{eq:#1}#2\end{equation}}
\newcommand{\Eqref}[1]{(\ref{eq:#1})}

% Common operators
\DeclareMathOperator*{\argmax}{arg\,max}
\DeclareMathOperator*{\argmin}{arg\,min}

% Probability / expectation / variance
\newcommand{\Prob}{\mathbb{P}}
\newcommand{\E}{\mathbb{E}}
\newcommand{\Var}{\mathrm{Var}}
\newcommand{\Cov}{\mathrm{Cov}}

% Sets, reals, indicators
\newcommand{\R}{\mathbb{R}}
\newcommand{\1}[1]{\mathbb{I}\{#1\}}

% Vectors / matrices
\newcommand{\vct}[1]{\bm{#1}}     % bold vector
\newcommand{\mtx}[1]{\bm{#1}}     % bold matrix
\newcommand{\trans}{^\mathsf{T}}  % transpose
\newcommand{\inv}{^{-1}}

% Norms, inner products
\newcommand{\norm}[1]{\left\lVert #1 \right\rVert}
\newcommand{\ip}[2]{\left\langle #1, #2 \right\rangle}

% Derivatives
\newcommand{\dd}[2]{\frac{d #1}{d #2}}
\newcommand{\pp}[2]{\frac{\partial #1}{\partial #2}}
\newcommand{\ppn}[3]{\frac{\partial^{#3} #1}{\partial #2^{#3}}}

% Convenience delimiters
\newcommand{\paren}[1]{\left( #1 \right)}
\newcommand{\brak}[1]{\left[ #1 \right]}
\newcommand{\set}[1]{\left\{ #1 \right\}}

% Theorem-like envs (optional)
\theoremstyle{plain}
\newtheorem{assumption}{Assumption}
\newtheorem{proposition}{Proposition}

% -----------------------
% Title block
% -----------------------
\begin{document}


\vspace{0.5em}

\begin{titlepage}
    \centering
    \vspace*{2cm}
    
\begin{center}
  {\large \textbf{\Title}}\\[0.5em]
  {\large \Author}\\
  {\Affil}\\[0.25em]
  {\Date}
\end{center}
    \vspace{2cm}
    
    \begin{minipage}{0.85\textwidth}
        \begin{center}
            \textbf{Abstract}
        \end{center}
        This paper investigates the impact of a women's professional sports league on girls' youth sports participation. Using the Women's National Basketball Association's (WNBA) inaugural season in 1997 as treatment, I apply two-way fixed effects methods to determine if there are any effects of the league's formation on the participation of girls in high school sports. Across four sports, I find modest increases in basketball and track participation, while cross-country and soccer do not exhibit any statistically significant effects.
    \end{minipage}
    
    \vfill
\end{titlepage}

\section{Introduction}

In 1997, the Women's National Basketball Association (WNBA) held its inaugural season. Since its conception, the league has expanded from an initial six teams, to a recently accounced 18 teams by 2030. This growth across states presents an interesting question: does the existence of a professional sports league encourage girls to engage in youth sports? 

\section{Literature}

\section{Methodology}

\subsection{Data}

\subsection{Estimation}

The proliferation of the WNBA lends itself to a differences-in-differences design. However, using the canonical differences-in-differences design would lead to biased estimates due to staggered treatment. Thus, I shift to the design of \cite{Callaway_SA}, which accounts for the bias by construction.


The estimating equation is as follows:

\begin{equation}
Y_{it} = \alpha_i + \lambda_t + \delta \cdot D_{it} + \varepsilon_{it}
\end{equation}

where $Y_{it}$ is state $i$'s participation in a sport at time $t$, $\alpha_i$ accounts for state-level fixed effects, $\lambda_t$ accounts for year effects, $D_{it}$ is a treatment indicator, indicating if state $i$ has a WNBA team in year $t$. Thus, $\delta$ represents the average treatment effect for a unit treated in a given expansion cohort. 



\section{Results}

\section{Robustness}

\section{Discussion}

\section{Conclusion}

\printbibliography
  

\end{document}
